% !TEX program = xelatex
% 非官方：仅据药学院2019规范，不代表全校本科统一格式。
\documentclass[12pt,a4paper,fontset=none]{ctexart}
\usepackage[top=25.4mm,bottom=25.4mm,left=31.7mm,right=31.7mm,headheight=15pt,headsep=5mm]{geometry}
% MIT; default preview uses TeX-distributed fonts, NOT school typography.
\RequirePackage{fontspec,xeCJK,etoolbox}
\providecommand{\ThesisFontMode}{preview}
\ifdefstring{\ThesisFontMode}{school}{%
  \newcommand{\ThesisCheckFont}[1]{\IfFontExistsTF{#1}{}{\PackageError{thesis-fonts}{Missing font: #1}{Install legally licensed fonts or use THESIS_FONTS=preview.}}}
  \ThesisCheckFont{Times New Roman}\ThesisCheckFont{SimSun}
  \ThesisCheckFont{SimHei}\ThesisCheckFont{FangSong}\ThesisCheckFont{KaiTi}
  \setmainfont{Times New Roman}
  \setCJKmainfont[AutoFakeBold=2]{SimSun}
  \setCJKsansfont[AutoFakeBold=2]{SimHei}
  \setCJKmonofont{FangSong}
  \setCJKfamilyfont{zhsong}[AutoFakeBold=2]{SimSun}
  \setCJKfamilyfont{zhhei}[AutoFakeBold=2]{SimHei}
  \setCJKfamilyfont{zhkai}[AutoFakeBold=2]{KaiTi}
  \setCJKfamilyfont{zhfs}{FangSong}
}{%
  \ifdefstring{\ThesisFontMode}{preview}{}{\PackageError{thesis-fonts}{Invalid font mode}{Use preview or school.}}
  \setmainfont{texgyretermes-regular.otf}[BoldFont=texgyretermes-bold.otf,ItalicFont=texgyretermes-italic.otf,BoldItalicFont=texgyretermes-bolditalic.otf]
  \setCJKmainfont{FandolSong-Regular.otf}[BoldFont=FandolSong-Bold.otf]
  \setCJKsansfont{FandolHei-Regular.otf}[BoldFont=FandolHei-Bold.otf]
  \setCJKmonofont{FandolFang-Regular.otf}
  \setCJKfamilyfont{zhsong}{FandolSong-Regular.otf}[BoldFont=FandolSong-Bold.otf]
  \setCJKfamilyfont{zhhei}{FandolHei-Regular.otf}[BoldFont=FandolHei-Bold.otf]
  \setCJKfamilyfont{zhkai}{FandolKai-Regular.otf}[AutoFakeBold=2]
  \setCJKfamilyfont{zhfs}{FandolFang-Regular.otf}
}
% Geometric/shapes may not exist in the serif font; ctex's open symbol font.
\xeCJKDeclareCharClass{CJK}{"2460 -> "24FF,"3200 -> "32FF}

\providecommand{\songti}{\CJKfamily{zhsong}}
\providecommand{\heiti}{\CJKfamily{zhhei}}
\xeCJKDeclareCharClass{CJK}{"2460 -> "24FF}
\usepackage{amsmath,amssymb,graphicx,booktabs,tabularx,longtable,caption}
\usepackage{titlesec,tocloft,fancyhdr,setspace,indentfirst,perpage,listings}
\usepackage{tikz,pgfplots}
\usetikzlibrary{positioning,arrows.meta}
\pgfplotsset{compat=1.18}
\usepackage{gbt7714}
\usepackage[hidelinks]{hyperref}
\setcounter{tocdepth}{3}\setcounter{secnumdepth}{3}
\titleformat{\section}{\centering\heiti\zihao{3}}{\thesection}{1em}{}
\titleformat{\subsection}{\heiti\zihao{-4}}{\thesubsection}{1em}{}
\titleformat{\subsubsection}{\heiti\zihao{-4}}{\thesubsubsection}{1em}{}
\renewcommand{\cftsecfont}{\heiti\zihao{-4}}
\renewcommand{\cftsubsecfont}{\songti\zihao{-4}}
\renewcommand{\cftsubsubsecfont}{\songti\zihao{-4}}
\renewcommand{\cftsecleader}{\cftdotfill{\cftdotsep}}
\renewcommand{\cfttoctitlefont}{\hfill\heiti\zihao{3}}
\renewcommand{\cftaftertoctitle}{\hfill\mbox{}}
\tocloftpagestyle{empty}
\fancypagestyle{medical}{\fancyhf{}\fancyhead[C]{\songti\zihao{5}海南医科大学本科毕业论文}\fancyhead[R]{\songti\zihao{5}第\thepage 页}\renewcommand{\headrulewidth}{0.4pt}}
\DeclareCaptionFont{medcap}{\heiti\zihao{5}}
\captionsetup{font=medcap,labelsep=space,justification=centering}
\captionsetup[table]{position=top}
\MakePerPage{footnote}
\renewcommand{\thefootnote}{\ifcase\value{footnote}\or①\or②\or③\or④\or⑤\or⑥\or⑦\or⑧\or⑨\or⑩\or⑪\or⑫\or⑬\or⑭\or⑮\or⑯\or⑰\or⑱\or⑲\or⑳\else\arabic{footnote}\fi}
\renewcommand{\footnotesize}{\zihao{-5}}
\lstset{basicstyle=\ttfamily\small,breaklines=true,columns=fullflexible,frame=single,numbers=left,numberstyle=\tiny,showstringspaces=false}

\newcommand{\thesistitle}{药学研究论文写作与排版示例}
\newcommand{\thesistitleeng}{A Typesetting Example for a Pharmaceutical Research Thesis}
\newcommand{\thesisauthor}{张三}
\newcommand{\thesisauthoreng}{Zhang San}
\newcommand{\thesisid}{2022000001}
\newcommand{\thesismajor}{药学}
\newcommand{\thesisadvisor}{李四}
\newcommand{\thesisadvisoreng}{Li Si}
\newcommand{\thesisdate}{二〇二六年六月}
\begin{document}
\songti\zihao{-4}\setstretch{1.5}\setlength{\parindent}{2em}
\begin{titlepage}\thispagestyle{empty}\centering
\vspace*{10mm}{\heiti\zihao{1}海南医科大学}\par\vspace{15mm}
{\heiti\zihao{2}本科毕业论文}\par\vspace{22mm}
{\heiti\zihao{2}\thesistitle}\par\vspace{20mm}
{\songti\zihao{-3}\renewcommand{\arraystretch}{1.8}
\begin{tabular}{r>{\raggedright\arraybackslash}p{.63\linewidth}}
学\quad 院：&药学院\\专\quad 业：&\thesismajor\\
姓\quad 名：&\thesisauthor\\学\quad 号：&\thesisid\\
指导教师：&\thesisadvisor\\完成日期：&\thesisdate
\end{tabular}}\vfill
{\zihao{5}示例封面：学院统一封面原件尚未取得，提交前必须替换或核对。}
\end{titlepage}
\pagestyle{empty}\section*{原创性声明（示例文本）}
本人声明：所提交的论文应为本人在指导教师指导下独立完成的研究成果。引用他人成果应明确注明来源。此处仅为声明占位文本，不是已取得的学院官方声明；提交时须替换为学院当年规定的原文，不得代签或预填签名。
\vspace{25mm}\begin{flushright}作者签名：\underline{\makebox[30mm]{}}\par\medskip
日期：\underline{\makebox[40mm]{}}\end{flushright}
\clearpage
\begin{center}{\heiti\zihao{-3}\thesistitle}\par\medskip{\heiti\zihao{3}摘\quad 要}\end{center}
\noindent\textbf{目的：}展示药学本科研究论文的结构式摘要与主要排版元素，不报告真实研究成果。
\textbf{方法：}依据药学院2019年公开规范整理页面、标题、图表、公式、脚注和参考文献样式，配置可复现的编译流程。
\textbf{结果：}提供封面占位、声明、中英文摘要、目录、研究设计、结果报告、讨论、参考文献、附录及致谢示例；所有数值均为人为构造的排版数据，不具有医学或药学证据意义。
\textbf{结论：}本文件可作为论文写作起点，但不等于学校审定模板。使用者须核对所在学院当年要求，用真实研究内容和经审查的材料替换示例，不能将模板文字作为本人研究提交。
\par\medskip\noindent{\heiti\zihao{3}关键字：}{\zihao{-4}药学；论文写作；结构式摘要；排版示例}
\clearpage
\begin{center}{\bfseries\zihao{-3}\thesistitleeng}\par\medskip
{\bfseries\zihao{4}Author: \thesisauthoreng\par Supervisor: \thesisadvisoreng}\par\medskip
\textbf{Abstract}\end{center}
\noindent\textbf{Objective:} To demonstrate the structure and typesetting of an undergraduate pharmaceutical thesis, without reporting empirical findings.
\textbf{Methods:} The publicly available 2019 School of Pharmacy requirements were used to configure page layout, headings, tables, figures, equations, notes, and references.
\textbf{Results:} This example provides front matter, a research-design section, synthetic data displays, a discussion, and supplementary material. No clinical or pharmaceutical inference can be drawn from the example data.
\textbf{Conclusion:} This unofficial template requires verification against current college instructions. All example material must be replaced with independently conducted and appropriately approved research.
\par\medskip\noindent\textbf{Keywords:} pharmacy; thesis writing; structured abstract; typesetting
\clearpage\tableofcontents
\clearpage\pagenumbering{arabic}\pagestyle{medical}
\section*{前言}\addcontentsline{toc}{section}{前言}
研究论文应清楚说明研究背景、待解决问题和研究目的。本模板以分析方法评价为叙述线索，演示医学与药学论文常见的材料与方法、结果、讨论及结论结构；没有开展实验，不作临床、药效或安全性判断。正式选题必须符合专业培养要求，不能用综述直接替代药学院规范所要求的研究性论文。

文献检索宜追溯到研究原文、正式指南或标准。此处引用真实的方法学资料仅作引文排版演示\cite{ichq2,jcgm2008}，不表示本研究已经完成方法验证。模板不得替代伦理审批、实验室安全审查、数据授权或研究注册。
\section{材料与方法}
\subsection{研究设计与适用范围}
说明研究类型、地点、时间范围、研究对象和分析单位。涉及人体、动物或个人健康数据时，应在开始研究前完成适用的伦理审查和知情同意，并在文中准确报告批准机构和编号。本例不提供任何虚构的批准编号。
\subsubsection{材料与仪器}
按实际使用记录试剂名称、等级、批号、来源及仪器型号。可复现性依赖具体条件，而不是仪器品牌的简单罗列。药品名称优先使用规范通用名，首次出现的缩写需给出全称。
\subsection{技术路线与质量控制}
图~\ref{fig:workflow} 展示方法研究的组织方式。真实流程还需报告样品处理、盲法、随机化、重复测定及质量控制规则，并区分技术重复与生物学重复。
\begin{figure}[htbp]\centering
\begin{tikzpicture}[node distance=6mm,box/.style={draw,minimum width=60mm,minimum height=9mm,align=center,font=\small},>=Stealth]
\node[box](a){明确研究问题与预设分析方案};
\node[box,below=of a](b){材料准备与质量控制};
\node[box,below=of b](c){测量记录与异常数据审核};
\node[box,below=of c](d){统计分析与结果解释};
\draw[->](a)--(b);\draw[->](b)--(c);\draw[->](c)--(d);
\end{tikzpicture}
\caption{研究流程示例（非实际实验流程）}\label{fig:workflow}
\end{figure}
\subsection{统计分析与误差表达}
预先明确主要指标、统计方法和适用前提。相关不等于因果，统计显著也不等于临床或药学意义。对于重复测量，均值与样本标准差可按式~\eqref{eq:sd} 计算：
\begin{equation}\label{eq:sd}
\bar{x}=\frac{1}{n}\sum_{i=1}^{n}x_i,\qquad s=\sqrt{\frac{\sum_{i=1}^{n}(x_i-\bar{x})^2}{n-1}}.
\end{equation}
式中，$n$ 为独立观测数，不能把同一样品的重复读数任意当作独立样本。本脚注演示每页重新编号\footnote{本文件的数值均为合成数据，不能作为任何药品质量、诊断或治疗决策依据。}。
\subsection{数据保存与分析复现}
记录原始文件、变量定义、处理脚本、软件版本和排除标准；研究完成后按机构规定归档。不得删除不符合预期的结果，也不应在看到数据后悄悄改变主要指标。一般统计报告可参考 SAMPL 等真实指南\cite{sampl2015}，随机试验、观察性研究、诊断研究分别有不同报告要求\cite{consort2010,strobe2007}。
\clearpage\section{结果}
\subsection{描述性结果的呈现}
表~\ref{tab:synthetic-med} 仅演示三线表、单位和小数位的编排，不报告真实含量测定结果。正文需要先引出表格，说明观测量和记录条件，避免只写“结果见表”。
\begin{table}[htbp]\centering\caption{合成数据示例（非药学实验结果）}\label{tab:synthetic-med}
{\zihao{5}\begin{tabular}{rrr}\toprule 标识 & 参考值/(mg/L) & 示例读数/(mg/L)\\\midrule
1&1.00&1.02\\2&2.00&1.98\\3&3.00&3.03\\4&4.00&3.99\\\bottomrule\end{tabular}}
\end{table}
\subsection{图形与不确定性}
图~\ref{fig:synthetic-med} 为合成数据散点示例。正式图形应标注坐标、单位、样本数及误差线的含义；不能将标准差、标准误和置信区间混用。只有四个人工数据点不能支持线性范围、检出限或准确度结论。
\begin{figure}[htbp]\centering\begin{tikzpicture}
\begin{axis}[width=.85\linewidth,height=55mm,xlabel={参考值/(mg/L)},ylabel={示例读数/(mg/L)},grid=major]
\addplot[only marks,black]coordinates{(1,1.02)(2,1.98)(3,3.03)(4,3.99)};
\end{axis}\end{tikzpicture}
\caption{合成数据散点图（非方法验证证据）}\label{fig:synthetic-med}
\end{figure}
\clearpage\section{讨论}
\subsection{结果解释与既有证据的关系}
讨论应回答研究问题，解释观察结果与已有证据的一致或不一致之处，明确何者为直接证据、何者为推测。不能根据单一指标推断药效、安全性或临床价值，也不能把仪器读数直接等同于生物学机制。
\subsection{局限与适用条件}
报告样本来源、测量偏倚、混杂因素和外推范围。没有真实试验时，应明确不存在可报告的实证结论，而不是编造“差异具有统计学意义”。本模板中的统计公式和图表只验证排版能力。
\section{结论}
正式结论应与研究目的逐项对应，使用具体、可核验、有边界的语言。模板不提供药学研究发现；请替换为实际研究得出的结果，并说明仍需进一步验证的内容。

\clearpage\begingroup
\renewcommand{\bibsection}{\section*{参考文献}\addcontentsline{toc}{section}{参考文献}}
\renewcommand{\bibfont}{\songti\zihao{-4}}
\bibliographystyle{gbt7714-2015-numeric}\bibliography{refs}
\endgroup
\clearpage\section*{附录}\addcontentsline{toc}{section}{附录}
\subsection*{附录 A\quad 变量与缩写说明}
\begin{longtable}{lll}\toprule 名称 & 英文 & 中文说明\\\midrule\endhead
SD & Standard deviation & 标准差\\
CI & Confidence interval & 置信区间\\
sample\_id & Sample identifier & 样本标识\\\bottomrule
\end{longtable}
\subsection*{附录 B\quad 最小计算程序}
\begin{lstlisting}[language=Python]
from statistics import mean, stdev
synthetic_values = [1.02, 0.99, 1.01]
print(mean(synthetic_values), stdev(synthetic_values))
\end{lstlisting}
实际附录可加入经批准的研究表单、数据字典、分析程序或详细推导。不要把含姓名、病案号等可识别信息的原始材料直接放入公开论文。

\clearpage\section*{致\quad 谢}\addcontentsline{toc}{section}{致谢}
药学院2019年规范要求致谢为200至500字，单独成页并放在参考文献之后。本页提供写作提示，不伪造作者的个人经历。正式致谢可说明指导教师在研究设计、方法选择、实验实施及论文修改中的具体帮助，感谢确实提供设备、样品、技术支持和数据协助的机构与人员，并按资助方要求准确填写项目名称和编号。感谢研究参与者时应保护隐私，不公开未经同意的个人信息。对于不符合作者资格但提供协助者，应在征得同意后说明其贡献；不得将没有参与研究的人员写入，也不应以致谢代替必要的署名、引用或利益冲突披露。请根据真实情况独立撰写并删除本段说明。
\end{document}
